\chapter{Fonctions : définir, représenter, modéliser}
\section*{Exercice 12.1 -- Images et antécédents}
Pour $f(x)=(x-1)^2$ sur $\R$, calculer $f(-2)$ et donner les antécédents de $4$, de $0$, puis de $-1$.

\section*{Exercice 12.2 -- Domaine}
Déterminer le domaine de $f(x)=\sqrt{3-x}/(x+2)$.

\section*{Exercice 12.3 -- Appartenance}
Le point $A(2,5)$ appartient-il à la courbe $y=x^2+1$ ? Et $B(-2,-3)$ ?

\section*{Exercice 12.4 -- Tableau insuffisant}
Deux fonctions prennent les valeurs $0,0,0$ en $-1,0,1$. Cela prouve-t-il qu'elles sont identiques ? Donner un contre-exemple.

\section*{Exercice 12.5 -- Modèle géométrique}
Une clôture de $30$ mètres entoure un rectangle. Si un côté vaut $x$, exprimer l'aire et donner le domaine des rectangles non dégénérés.

\section*{Exercice 12.6 -- Modèle discret}
Un envoi de $n$ colis coûte $u(n)=8+2{,}5n$, pour $n\in\N$. Calculer $u(4)$ et expliquer si le prix $19$ euros est réalisable.
